Tais funções são (ibid.)

\begin{enumerate}
\item 🚜 $\rightarrow$ 🏭 mover materiais,
\item 🔒 liga/desliga máquinas,
\item 🧭 controlar volume de prod.,
\item 🕸️ ``monitorar/supervisionar o processo todo'' [\hl{IA} deseja \hl{automatizar} isso!!],
\item ``Handling of non automatic machines and tools'' $\dots$ não é sua função padrão...?
\end{enumerate}

\textbf{``Math theory of comm.'' de Shannon é de 1948.:} 
\begin{quote}
``The aims and methods of automation may be \textbf{provisionally} defined as a \textit{technique of production} the object of which is \textit{to replace men by machines} in \textit{operating and directing machines}[,] as well as in \textit{controlling the output} of the products that are being manufactured. If this aspect of automation is complete[,] the product is \textit{not touched} by human hands from the beginning to the end of its manufacture. Automation may be applied to \textit{part} of the process of manufacture or to the \textit{entire} process of turning raw materials into finished products.''  \parencite[p. 5, grifo meu]{Pollock1957}
\end{quote}

\begin{itemize}
\item a produção de $N$ mercadorias requer, como na média, o mesmo adiantamento de capital $c+v$, porém é produzida em $\Delta t^{\prime}<\Delta t$;
\item a produção de $\Delta N$ mercadorias, no período $\Delta t-\Delta t^{\prime}>0$, requer um adiantamento a mais de Capital Constante, \textit{porém emprega o mesmo capital variável adiantado $v$}, pois ele foi contratado pelo período de tempo $\Delta t$ \textit{inteiro}, não somente $\Delta t^{\prime}$.
\item Portanto, no período de tempo $\Delta t$ inteiro, o capital adiantado para este capital mais produtivo será $(1+\sigma)c + v$, produzindo $N+\Delta N$ mercadorias, e valor (individual) total $(1+\sigma)c + (1+m^{\prime})v$.
\end{itemize}

Para que compense para este capital produzir a este nível maior de produtividade, é necessário que seu ``custo'' (valor ``individual'') seja menor (ou, no máximo, igual) que o valor social:
\begin{align*}
(\cancel{ 1 }+\sigma)c + \cancel{ (1 +m^{\prime})v } &\leq \cancel{ c }+\cancel{ (1+m^{\prime})v } + \frac{\Delta N}{N}(c+(1+m^{\prime})v) \\
\iff \sigma c & \leq \frac{\Delta N}{N}[c+(1+m^{\prime})v]
\end{align*}

Dito em termos ``vulgares'': o ``custo extra'' ($\sigma c$) tem que ser menor (ou igual) ao ``lucro extra'' $\frac{\Delta N}{N}(c+v)$ de $\Delta N$ mercadorias a mais sendo vendidas ao valor social unitário $\frac{c+v}{N}$. \footnote{Test test \textbf{\textit{one two three}}.}

\begin{cases}
x&: \text{Ativo arriscado}\\
w-x&: \text{Riqueza não-investida} 
\end{cases}

Mostre que a solução ótima de riqueza pós-taxação $\hat{x}$ de se investir satisfaz\footnote{Sinta-se livre para explorar este exercício nesta simulação: \href{https://www.desmos.com/calculator/ppkg28wznf}{Investimentos e Taxação (Simulação no Desmos)}!}
$$
\hat{x} = \frac{x^*}{1-t}
$$
quando $x^\textit{$ é solução ótima pré-taxação. Isso condiz com sua intuição de $\hat{x} \gtrless x^}$?
